\section{Triton mental model}

Source coverage: \texttt{triton\_introduction / triton\_gelu\_example}. 本节覆盖的核心概念包括：Triton program, block, mask, pointer offsets, PTX。

\begin{knowledgebox}{课堂提示：Triton 把思考单位从 thread 提升到 block}
老师对比 CUDA 与 Triton：CUDA 让程序员细粒度指定每个 thread 做什么，控制强但也要手动管理更多 shared-memory 与调度细节；Triton 让程序员描述一个 thread block/program 如何加载一块数据、在片上操作、再写回 global memory。对深度学习 kernel，这个抽象通常已经足够强，同时更接近 tensor 运算的数学结构。
\end{knowledgebox}


\begin{figure}[H]
\centering
\includegraphics[width=0.86\textwidth]{cuda-grid.png}
\caption{cuda grid.png}
\figsource{CS336 Spring 2026 Lecture 6 官方可执行讲义。}
\end{figure}
\begin{knowledgebox}{读图：cuda grid.png}
这张图用于把抽象的 kernel 代码连接到硬件执行。读图时先看数据位于 HBM、shared memory 还是 registers，再看线程/blocks 如何映射到 SM。性能优化的关键是让数据在更靠近计算单元的位置被复用，减少慢速 HBM 往返。
\end{knowledgebox}

\begin{lstlisting}[caption={Triton block-level 思维}]
pid = tl.program_id(axis=0)
offsets = pid * BLOCK_SIZE + tl.arange(0, BLOCK_SIZE)
mask = offsets < num_elements
x = tl.load(x_ptr + offsets, mask=mask)
tl.store(y_ptr + offsets, y, mask=mask)
\end{lstlisting}
\begin{knowledgebox}{Triton 与 CUDA 的差异}
CUDA 通常让你描述每个 thread 做什么；Triton 更像描述一个 program/block 如何处理一块数据。它牺牲部分底层控制，换来更高层、更适合深度学习 kernel 的表达。
\end{knowledgebox}
\subsection{本章小结}
本节说明了一个 kernel optimization 的共同模式：先确认语义正确，再测时间，再看 profiler，最后用 block、shared memory、tiling、fusion 或 layout 调整降低数据移动。

